% --------------------------------------------------------------------------
\section{Complete Training Framework}
\label{sec:ch4_complete_pipeline}

AttackDRO++ is the final method studied in this report. It combines latent
cluster discovery, augmented clustering features, gradient fingerprints, and a
uniform-anchored Cluster-DRO objective into one training pipeline. The goal is to
move beyond fixed attack identities and let the model identify harder adversarial
subgroups during training. At the same time, the uniform anchor prevents the
adaptive objective from drifting too far from the strong multi-attack baseline.
Algorithm~\ref{alg:ch4_attackdropp} gives the complete procedure, while
Table~\ref{tab:ch4_hparams} lists the default hyperparameters used in the main
experiments.

\paragraph{From Diagnostic Motivation to the Final Method}

The method follows directly from the attacks-as-domains diagnosis introduced
earlier. Single-attack adversarial training reveals the cross-attack robustness
gap: a model trained against one source attack can remain weak against other
threat models. Uniform multi-attack training addresses this by exposing the model
to multiple source attacks, but it treats all generated examples equally.
AttackDRO then adds attack-level reweighting, but this still assumes that each
attack family is internally homogeneous.

AttackDRO++ removes this assumption. Instead of using attack identity as the only
group label, it discovers latent groups from the adversarial examples themselves.
Gradient fingerprints add optimization-level information to the clustering
features, helping the grouping capture not only representation similarity but also
training difficulty. Finally, the uniform anchor keeps the objective close to the
strong multi-attack baseline while still allowing Cluster-DRO to emphasize harder
groups. The result is a method that moves from observing the cross-attack gap to
actively training against the latent structure behind it.

The same evaluation discipline is preserved throughout the framework. Source
attacks are fixed during training, held-out attacks are never used for
optimization, and robustness is reported both per attack and through aggregate
metrics.

The dominant cost remains adversarial example generation, which is shared by all
multi-attack methods in the comparison. Therefore, AttackDRO++ mainly changes the
training objective and grouping mechanism, while keeping the source attacks,
threat models, and evaluation protocol fixed. The cluster refresh schedule is an implementation choice rather than a fixed part
of the mathematical objective. In the main experiments, clusters are refreshed
every \(R=2\) epochs, and \(\mathbf{q}\) is reset to the uniform distribution
after each refresh to avoid carrying weights across changed cluster identities.


\begin{figure}[htbp]
\centering
\resizebox{0.95\textwidth}{!}{%
\begin{tikzpicture}[
  node distance = 0.68cm,
  every node/.append style = {align=center},
  inputblk/.style  = {draw=black!40,        fill=gray!11,
                      rounded corners=5pt,  minimum width=7.6cm,
                      minimum height=1.08cm, font=\small,
                      inner sep=8pt,        line width=0.85pt},
  attackblk/.style = {draw=red!50!black,    fill=red!9,
                      rounded corners=5pt,  minimum width=7.6cm,
                      minimum height=1.08cm, font=\small,
                      inner sep=8pt,        line width=0.85pt},
  fwdblk/.style    = {draw=blue!55!black,   fill=blue!8,
                      rounded corners=5pt,  minimum width=7.6cm,
                      minimum height=1.08cm, font=\small,
                      inner sep=8pt,        line width=0.85pt},
  featblk/.style   = {draw=violet!65!black, fill=violet!8,
                      rounded corners=5pt,  minimum width=7.6cm,
                      minimum height=1.08cm, font=\small,
                      inner sep=8pt,        line width=0.85pt},
  clustblk/.style  = {draw=orange!75!black, fill=orange!10,
                      rounded corners=5pt,  minimum width=7.6cm,
                      minimum height=1.08cm, font=\small,
                      inner sep=8pt,        line width=0.85pt},
  lossblk/.style   = {draw=green!60!black,  fill=green!9,
                      rounded corners=5pt,  minimum width=3.5cm,
                      minimum height=1.60cm, font=\small,
                      inner sep=8pt,        line width=0.85pt},
  combblk/.style   = {draw=teal!65!black,   fill=teal!9,
                      rounded corners=5pt,  minimum width=7.6cm,
                      minimum height=1.08cm, font=\small,
                      inner sep=8pt,        line width=0.85pt},
  updblk/.style    = {draw=teal!80!black,   fill=teal!13,
                      rounded corners=5pt,  minimum width=7.6cm,
                      minimum height=1.08cm, font=\small,
                      inner sep=8pt,        line width=0.85pt},
  sideblk/.style   = {draw=gray!55,         fill=gray!8,
                      rounded corners=5pt,  minimum width=3.1cm,
                      minimum height=1.05cm, font=\small,
                      inner sep=7pt,        line width=0.72pt,
                      densely dashed},
  arr/.style       = {-{Stealth[length=7pt, width=5.5pt]},
                      line width=0.95pt, color=black!55},
  darr/.style      = {-{Stealth[length=6pt, width=4.5pt]},
                      line width=0.72pt, color=gray!60, densely dashed},
  looparr/.style   = {-{Stealth[length=7pt, width=5.5pt]},
                      line width=0.80pt, color=black!28, rounded corners=4pt},
]

\node[inputblk] (batch)
  {\textbf{Minibatch}\qquad
   $\mathcal{B} = \{(x_i,\,y_i)\}_{i=1}^{N}$};

\node[attackblk, below=of batch] (attacks)
  {\textbf{Mixed Source-Attack Batch}\\[3pt]
   split minibatch across
   $\mathrm{PGD}\text{-}\ell_\infty$ and $\mathrm{DDN}\text{-}\ell_2$
   \quad$\longrightarrow$\quad
   $\{x_i^{\mathrm{adv}}\}_{i=1}^{N}$};

\node[fwdblk, below=of attacks] (forward)
  {\textbf{Model Signals from } $f_\theta(x_i^{\mathrm{adv}})$\\[3pt]
   forward: logits $z_i$, representation $h_i$
   \quad$|$\quad
   gradient: fingerprint $g_i$};

\node[featblk, below=of forward] (augfeat)
  {\textbf{Augmented Clustering Feature}\\[3pt]
   $\psi_i = \bigl[\,h_i \;;\; z_i \;;\; y_i \;;\; g_i\,\bigr]$};

\node[clustblk, below=of augfeat] (cluster)
  {\textbf{Cluster Assignment}\enspace (current $k$-means)\\[3pt]
   $c_i \;\in\; \{1,\dots,K\}$};

\node[lossblk] (lerm) at ($(cluster.south)+(-2.8,-1.95)$)
  {$\mathcal{L}_{\mathrm{ERM}}$\\[6pt]
   {\footnotesize uniform multi-attack}};

\node[lossblk] (ldro) at ($(cluster.south)+(+2.8,-1.95)$)
  {$\mathcal{L}_{\mathrm{ClusterDRO}}$\\[6pt]
   $\displaystyle\sum_{k=1}^{K} q_k\,\mathcal{L}_k$};

\coordinate (lossmid) at ($(lerm.south)!0.5!(ldro.south)$);

\node[combblk] (combine) at ($(lossmid)+(0,-1.35)$)
  {$\mathcal{L}_{\mathrm{total}}
    \;=\;
    (1-\lambda_{\mathrm{DRO}})\,\mathcal{L}_{\mathrm{ERM}}
    \;+\;
    \lambda_{\mathrm{DRO}}\,\mathcal{L}_{\mathrm{ClusterDRO}}$};

\node[updblk, below=0.65cm of combine] (updatetheta)
  {\textbf{Update Model Parameters}\\[3pt]
   $\theta \;\leftarrow\; \theta
      - \eta_\theta\,\nabla_\theta\,\mathcal{L}_{\mathrm{total}}$};

\node[updblk, below=0.65cm of updatetheta] (updateq)
  {\textbf{Update Cluster Weights}\enspace
   (EG update $+$ normalisation $+$ $q$\text{-floor})\\[3pt]
   $q_k \;\leftarrow\;
     \mathrm{Proj}_{q_{\min}}\!\!\left(
       q_k \cdot \exp\!\bigl(\eta_q\,\mathcal{L}_k\bigr)\right)$};

\node[sideblk] (bank) at ($(augfeat.east)+(3.75,0)$)
  {\textbf{Feature Bank}\\[3pt]
   $\{\psi_i\}$ \;(detached)};

\node[sideblk, below=1.50cm of bank] (kmeans)
  {\textbf{$k$-Means Clusterer}\\[3pt]
   $K$ centroids};

\draw[arr] (batch)     -- (attacks);
\draw[arr] (attacks)   -- (forward);
\draw[arr] (forward)   -- (augfeat);
\draw[arr] (augfeat)   -- (cluster);

\coordinate (splitpt) at ($(cluster.south)+(0,-0.44)$);
\draw[arr] (cluster.south) -- (splitpt);
\draw[arr] (splitpt) -| (lerm.north);
\draw[arr] (splitpt) -| (ldro.north);

% Merge the two loss branches before entering the total-loss box.
\coordinate (lossJoinR) at ($(ldro.south)+(0,-0.32)$);
\coordinate (lossJoinL) at (lerm.south |- lossJoinR);
\coordinate (lossJoinM) at ($(lossJoinL)!0.5!(lossJoinR)$);
\draw[line width=0.95pt, color=black!55] (lerm.south) -- (lossJoinL);
\draw[line width=0.95pt, color=black!55] (ldro.south) -- (lossJoinR);
\draw[line width=0.95pt, color=black!55] (lossJoinL) -- (lossJoinR);
\draw[arr] (lossJoinM) -- (combine.north);

\draw[arr] (combine)     -- (updatetheta);
\draw[arr] (updatetheta) -- (updateq);

\draw[darr] (augfeat.east)
    to node[above, font=\scriptsize\itshape, color=gray!65]
       {store (detach)}
    (bank.west);

\draw[darr] (bank.south)
    to node[right, font=\scriptsize\itshape, color=gray!65,
            xshift=2pt, align=left]
       {periodic\\[-2pt]refresh\\[-2pt]($R$ epochs)}
    (kmeans.north);

\draw[darr] (kmeans.west) -- ++(-0.65,0)
    |- node[left, font=\scriptsize\itshape, color=gray!65,
            xshift=-2pt, near start] {assign}
    (cluster.east);

\coordinate (loopbot)  at ($(updateq.south)+(0,-0.58)$);
\coordinate (loopleft) at ($(loopbot)+(-6.80,0)$);
\draw[looparr]
    (updateq.south) -- (loopbot) -- (loopleft) |- (batch.west);

\node[font=\scriptsize\itshape, color=black!32,
      anchor=north, inner sep=2pt]
    at ($(loopbot)!0.5!(loopleft)$)
    {next iteration};

\end{tikzpicture}%
}
\caption[Core AttackDRO++ training loop]{%
  \textbf{Core training loop of AttackDRO++.}
  At each iteration, the minibatch is split across two source attacks
  ($\mathrm{PGD}\text{-}\ell_\infty$ and $\mathrm{DDN}\text{-}\ell_2$),
  producing a mixed adversarial batch.
  The batch is passed through \(f_\theta\) to obtain logits \(z_i\),
  penultimate representations \(h_i\), and gradient fingerprints \(g_i\).
  These signals are concatenated into augmented features
  \(\psi_i=[h_i;z_i;y_i;g_i]\) and assigned to one of \(K\) latent clusters by
  the current \(k\)-means clusterer.
  The objective interpolates a uniform multi-attack ERM term
  \(\mathcal{L}_{\mathrm{ERM}}\) with a cluster-DRO term
  \(\mathcal{L}_{\mathrm{ClusterDRO}} = \sum_k q_k\mathcal{L}_k\) using
  anchor coefficient \(\lambda_{\mathrm{DRO}}\).
  Model parameters \(\theta\) are updated every step, while cluster weights
  \(\mathbf{q}\) are updated after the warmup period using exponentiated-gradient
  updates with floor projection; the detached feature bank periodically refreshes
  the \(k\)-means centroids every \(R\) epochs.%
}
\label{fig:attackdropp_core_training}
\end{figure}


\begin{algorithm}[H]
\caption{AttackDRO++ with Gradient Fingerprints and Uniform Anchor}
\label{alg:ch4_attackdropp}
\DontPrintSemicolon
\SetAlgoLined

\KwIn{Training set \(\mathcal{D}_{\mathrm{train}}\), model \(f_\theta\), source attacks \(\mathcal{A}_{\mathrm{src}}=\{A_1,\ldots,A_M\}\), Cluster count \(K\), anchor weight \(\lambda_{\mathrm{DRO}}\), learning rates \(\eta_\theta,\eta_q\), \(q\)-floor \(q_{\min}\), Label weight \(\lambda_{\mathrm{lbl}}\), gradient weight \(\lambda_{\mathrm{grad}}\), projection dimension \(d_{\mathrm{proj}}\), Refresh interval \(R\), warmup epochs \(T_w\)}
\KwOut{Trained model \(f_\theta^\ast\)}

Initialize \(\mathbf{q} \leftarrow (1/K,\ldots,1/K)\)\;
Initialize feature bank \(\mathcal{F} \leftarrow \emptyset\)\;
Initialize K-means clusterer as unavailable\;
Sample fixed projection matrix \(P \in \mathbb{R}^{C d_h \times d_{\mathrm{proj}}}\)\;

\For{epoch \(e = 1,\ldots,E\)}{
    \For{each minibatch \((x,y)\) of size \(N\)}{
        Generate mixed adversarial batch \(x^{\mathrm{adv}}\) using \(\mathcal{A}_{\mathrm{src}}\)\;
        Compute logits and features \((z,h) \leftarrow f_\theta(x^{\mathrm{adv}})\)\;
        Build augmented features \(\psi \leftarrow \mathrm{BuildFeature}(z,h,y,P)\)\;
        Assign cluster ids \(c \leftarrow \mathrm{AssignCluster}(\psi,\mathrm{Clusterer},K)\)\;

        Compute uniform loss \(\mathcal{L}_{\mathrm{ERM}}\)\;
        Compute cluster-DRO loss \(\mathcal{L}_{\mathrm{ClusterDRO}}\)\;
        Compute final loss
        \[
        \mathcal{L}
        =
        (1-\lambda_{\mathrm{DRO}})\mathcal{L}_{\mathrm{ERM}}
        +
        \lambda_{\mathrm{DRO}}\mathcal{L}_{\mathrm{ClusterDRO}} .
        \]

        Update model parameters \(\theta\) using SGD\;

        \If{\(e > T_w\)}{
            Update \(\mathbf{q}\) using exponentiated gradient and floor projection\;
        }

        Add detached features \(\psi\) to feature bank \(\mathcal{F}\)\;
    }

    \If{refresh condition is met}{
        Refresh K-means clusterer using \(\mathrm{KMeans}(\mathcal{F},K)\)\;
        Apply \(q\)-refresh policy to \(\mathbf{q}\) \tcp*{default: reset to uniform}
        Reset or update feature bank \(\mathcal{F}\) according to the memory policy\;
    }
}

\Return{\(f_\theta^\ast\)}
\end{algorithm}


\begin{table}[H]
\centering
\caption{Default hyperparameters for AttackDRO++ in the main experiments.}
\label{tab:ch4_hparams}
\small
\begin{tabularx}{\textwidth}{@{}l l l X@{}}
\toprule
Hyperparameter & Symbol & Value & Role \\
\midrule
DRO anchor strength      & \(\lambda_{\mathrm{DRO}}\)  & 0.35  & Interpolates between uniform ERM and Cluster-DRO \\
Number of clusters       & \(K\)                       & 4     & Controls latent group granularity \\
\(q\) learning rate      & \(\eta_q\)                  & \(3\times10^{-4}\) & Step size for exponentiated-gradient update \\
\(q\)-floor              & \(q_{\min}\)                & 0.03  & Minimum mass retained by each cluster \\
\(q\)-warmup             & \(T_w\)                     & 3 epochs & Delays adaptive reweighting until clusters are more reliable \\
Gradient projection dim. & \(d_{\mathrm{proj}}\)       & 128   & Dimension of projected gradient fingerprint \\
Gradient weight          & \(\lambda_{\mathrm{grad}}\) & 0.1   & Weight of GradFP component in \(\psi\) \\
Label weight             & \(\lambda_{\mathrm{lbl}}\)  & 0.1   & Weight of label component in \(\psi\) \\
Feature bank size        & \(|\mathcal{F}|\)           & 4{,}096 & Buffer used to refit K-means \\
Cluster refresh interval & \(R\)                       & 2 epochs & Default K-means refit frequency \\
\(q\)-refresh policy     & N/A                         & reset & Resets \(\mathbf{q}\) to uniform after cluster refresh \\
Bootstrap assignment     & N/A                         & random & Used before the first K-means refit \\
\bottomrule
\end{tabularx}
\end{table}
